\section{Model og tidsrespons}
\label{sec:model}

\subsection{ODE eller fysisk model til transferfunktion / Laplace, modelling}

\paragraph{Brug når}
Opgaven giver en differentialligning, et RLC-kredsløb, masser/fjedre eller beder om
poler og nulpunkter.

\paragraph{Input fra opgaven}
Input/output-definition, fysiske parametre, ODE-led og begyndelsesbetingelser.

\paragraph{Metode}
\begin{enumerate}
\item Vælg positiv retning for hvert signal og skriv den på figuren eller i teksten.
      For mekanik: tegn free-body diagram og brug \(\sum F=m\ddot x\) i den valgte
      positive retning. For kredsløb: vælg strømretning og spændingspolariteter, og brug
      KCL/KVL med samme fortegn i alle led.
\item Skriv én ligning pr. fysisk lov eller pr. opgivet ODE. Hvis opgaven allerede giver
      en lineær ODE, må du ikke kalde den ikke-lineær; gå direkte til Laplace-trinnet.
      Hvis et led er ikke-lineært, vælg arbejdspunkt \((x_0,u_0)\), skriv
      \(x=x_0+\tilde x\), \(u=u_0+\tilde u\), behold konstant- og førsteordensled fra
      Taylor/Jacobian, og drop produkter som \(\tilde x\tilde u\).
\item Sæt initialbetingelser til nul, medmindre opgaven udtrykkeligt spørger om dem.
      Erstat derefter hvert \(y^{(n)}(t)\) med \(s^nY(s)\), hvert \(u^{(n)}(t)\) med
      \(s^nU(s)\), og saml alle \(Y\)-led på venstre side og \(U\)-led på højre side.
\item Dan transferfunktionen ved at dividere med inputtet: \(G(s)=Y(s)/U(s)\) i
      \eqref{eq:transfer-function-dc}. Poler er rødder i nævneren efter fælles faktorer
      er forkortet; zeros er rødder i tælleren.
\item Klassificer stabilitet før DC-gain eller slutværdi: alle poler skal ligge strengt
      i \(\LHP\). Hvis der er en RHP-pol, en gentagen imaginæraksepol eller en origo-pol
      i den lukkede respons, må du ikke bruge en endelig slutværdi.
\end{enumerate}

\paragraph{Formler}
\begin{equation}
\label{eq:transfer-function-dc}
G(s)=\frac{Y(s)}{U(s)}=\frac{b_ms^m+\cdots+b_0}{a_ns^n+\cdots+a_0},
\qquad K_{\mathrm{ss}}=\lim_{s\to0}G(s).
\end{equation}
Alle poler strengt i \(\LHP\) giver asymptotisk stabil rational kontinuert model.
En RHP-pol eller gentagen pol på imaginæraksen er ustabil i kursets betydning.
\begin{equation}
\label{eq:first-final-value}
y(0^+)=\lim_{s\to\infty}sY(s),\qquad
y(\infty)=\lim_{s\to0}sY(s).
\end{equation}
Initial value theorem kræver at signalet ikke har impuls ved \(t=0\). Final value theorem
kræver at alle poler i \(sY(s)\), bortset fra en simpel pol i origo fra et stepinput,
ligger strengt i \(\LHP\).

\paragraph{Symboler}
\(Y,U\) er Laplace-transformer; \(p\) og \(z\) er nævner- og tællerrødder;
\(K_{\mathrm{ss}}\) er endelig DC-gain når den eksisterer.

\paragraph{Antagelser}
Transferfunktionmetoden kræver lineær tidsinvariant model; ODE-omskrivning ovenfor
forudsætter nul begyndelsesbetingelser, medmindre initialled håndteres eksplicit.

\paragraph{Hurtigtjek}
Nævnerens grad skal være mindst tællerens for en proper fysisk standardplant; kontrollér
origo-poler før du skriver en endelig DC-gain. Ved et input som \(u(t)=e^{-at}+\theta(t)\)
skal du først skrive \(U(s)=1/(s+a)+1/s\), derefter \(Y(s)=G(s)U(s)\), og først så bruge
\eqref{eq:first-final-value}.

\paragraph{Typiske fælder}
At kalde en lineær ODE ikke-lineær; at overse \(s^2\) i F21 Q12; at læse et diagram uden
at bevare fortegn.

\paragraph{Python}
\code{transfer\_function\_poles(denominator)} kontrollerer poler efter din udledning.
Angiv koefficienter i faldende potenser af \(s\). Funktionen erstatter ikke modelleringen
og må ikke bruges på et ufortolket diagram.
\code{analyze\_transfer\_function(G, variable=s)} bruges, når du allerede har en
symbolsk transferfunktion og vil have samlet output: poler, nulpunkter, DC-gain,
systemtype, stabilitet, tidsrespons-estimater og frekvensmetrics. Brug den ikke som
erstatning for at vælge det rigtige input/outputpar fra et diagram.
\code{frequency\_response\_point(num,den,omega)} bruges, når en kendt model skal
evalueres ved én frekvens, og du vil have både kompleks værdi, magnitude, dB og fase.
\code{evaluate\_transfer\_function(num,den,omega)} er den korte wrapper, når kun
den komplekse værdi \(G(j\omega)\) ønskes.

Hvis et realt gain \(K\) indgår i det karakteristiske polynomium \(p(s,K)\), bruges
\code{find\_stable\_gain\_ranges(p, K, s)}. Den accepterer også den rationale
karakteristiske ligning, fx \(1+K/(s+1)^3=0\), og samler automatisk brøken før
den indsætter \(s=j\omega\), finder marginale gainværdier og tester intervallerne. For
\(p=(s+1)^3+K\) fås stabilitet for \(-1<K<8\), eller \(0<K<8\) når positivt gain
kræves. Grænsepunkter med poler på imaginæraksen tælles ikke som asymptotisk stabile.

\subsection{White-, grey- og black-box modellering / system identification}
\label{sec:modelling-types}

\paragraph{Brug når}
Slides eller opgave spørger om modeltype, valg af transferfunktionsstruktur,
parameterestimat fra data eller REGBOT-transferfunktion fra målinger.

\paragraph{Metode}
\begin{enumerate}
\item White-box: start i fysiske love, angiv antagelser, lineariser om et arbejdspunkt
      og udled \(G(s)\). Parametre har fysisk betydning.
\item Grey-box: brug kendt fysisk struktur, men estimer ukendte parametre fra data.
      Kontroller at de estimerede værdier stadig er fysisk rimelige.
\item Black-box: vælg en modelstruktur, fx pol-/zero-orden og eventuel delay, og fit
      koefficienter direkte til input-output-data.
\item Brug små eksitationssignaler ved identifikation, hvis modellen skal være lineær.
      Store steps kan aktivere saturation, rate limit eller friktion og give en model,
      der ikke beskriver small-signal-dynamikken.
\item Valider modellen på andre data end fit-data. En lav fit-fejl er ikke nok, hvis
      poler, DC-gain, delay eller responsretning er fysisk forkert.
\end{enumerate}

\paragraph{Formler og tommelfingerregler}
Førsteordens black-box fra et step: \(K_{\mathrm{ss}}=\Delta y/\Delta u\), og
\(\tau\) aflæses hvor output har nået ca. \(63.2\%\) af den samlede ændring. Har
responsen et tydeligt dødtidsstykke, modelleres det separat som \(e^{-sL}\) eller
approksimeres kun hvis opgaven kræver en rational transferfunktion.

\paragraph{Hurtigtjek}
Modelstrukturen skal kunne forklare både fortegn på DC-gain, startretning, dominerende
tidskonstant og eventuel delay. REGBOT-data i slides bruges som et eksempel på, at samme
fysiske system kan behandles som white-box i balanceudledningen og black-box ved
målebaseret positions-/hastighedsmodel.

\paragraph{Typiske fælder}
At kalde en datatilpasset model ``white-box''; at bruge et fit fra store signaler til
linear controllerdesign; at glemme at en black-box model uden fysisk tolkning stadig
skal kontrolleres for stabilitet og properness.

\subsection{Linearisering og tilstandsform / operating point, Jacobian}
\label{sec:linearization}

\paragraph{Brug når}
Opgaven giver en ikke-lineær ligning, et inverted-pendulum/REGBOT-arbejdspunkt eller
beder om en lokal transferfunktion omkring et stationært punkt.

\paragraph{Metode}
\begin{enumerate}
\item Find arbejdspunktet ved at sætte tidsafledte lig nul og løse
      \(0=f(x_0,u_0)\). Arbejdspunktet skal passe til den driftssituation, controlleren
      skal bruges i.
\item Skriv deviation variables: \(\tilde x=x-x_0\), \(\tilde u=u-u_0\),
      \(\tilde y=y-y_0\). Det er disse små signaler, Laplace-transferfunktionen
      beskriver.
\item For \(\dot x=f(x,u)\), \(y=g(x,u)\) bruges Jacobian:
      \(A=\partial f/\partial x|_0\), \(B=\partial f/\partial u|_0\),
      \(C=\partial g/\partial x|_0\), \(D=\partial g/\partial u|_0\).
\item Den lokale model er \(\dot{\tilde x}=A\tilde x+B\tilde u\),
      \(\tilde y=C\tilde x+D\tilde u\). Transferfunktionen er
      \(G(s)=C(sI-A)^{-1}B+D\), når denne form er nødvendig.
\item Efter linearisering: kontroller polerne. For REGBOT balance giver tilt/velocity-
      planten typisk en ustabil lokal model, som skal stabiliseres før performance
      designes.
\end{enumerate}

\paragraph{Typiske fælder}
At blande absolutte signaler og deviationssignaler; at linearisere omkring et punkt,
der ikke er stationært; at bruge den lokale lineære model til store vinkler eller
store referencespring.

\subsection{Første- og andenordens steprespons / time response, overshoot}
\label{sec:time}

\paragraph{Brug når}
Der spørges om tidskonstant, steprespons, overshoot, rise/settling time eller et gain,
der skal overholde et overshootkrav.

\paragraph{Input fra opgaven}
Closed-loop-transferfunktion eller nævner, stepstørrelse og performancekrav.

\paragraph{Metode}
\begin{enumerate}
\item For første orden: skriv transferfunktionen som \(K/(\tau s+1)\). Hvis nævneren
      er \(as+b\), divider med \(b\), så formen bliver \((a/b)s+1\); dermed er
      \(\tau=a/b\), og DC-gain findes med \eqref{eq:transfer-function-dc}.
\item For anden orden: divider nævneren med koefficienten foran \(s^2\), så den bliver
      \(s^2+a_1s+a_0\). Match derefter mod nævneren i \eqref{eq:second-order-step}:
      \(\omega_n=\sqrt{a_0}\) og \(\zeta=a_1/(2\omega_n)\).
\item Beregn overshoot med \eqref{eq:second-order-step}. Hvis kravet er fx højst
      \(12\%\), brug \(M_p\leq0.12\), ikke \(12\). Hvis \(K\) indgår i \(a_0\) eller
      \(a_1\), indsæt udtrykket for \(\zeta(K)\) og løs uligheden.
\item Hvis opgaven spørger om settling inden for \(1\%\) eller et andet bånd, brug
      eksponentialreglen i \eqref{eq:settling-band}. For en reel dominerende pol
      \(p_{\mathrm{dom}}\) er \(\tau_{\mathrm{dom}}=-1/p_{\mathrm{dom}}\). For et
      underdæmpet polpar er envelope-tidskonstanten \(1/(\zeta\omega_n)\).
\item Kontrollér at standardformen må bruges. For ren første-/andenorden er den eksakt.
      For højere orden kræves at ekstra poler er mindst ca. 5 gange længere mod venstre
      end de dominerende poler, eller at opgaven eksplicit beder om en
      dominant-pole-approksimation.
\end{enumerate}

\paragraph{Formler}
\begin{equation}
\label{eq:first-order-step}
G_1(s)=\frac{K_{\mathrm{ss}}}{\tau s+1},\quad
y(t)=h_0K_{\mathrm{ss}}(1-e^{-t/\tau}),\quad \omega_b=\frac1\tau .
\end{equation}
\begin{equation}
\label{eq:second-order-step}
G_2(s)=\frac{\omega_n^2}{s^2+2\zeta\omega_ns+\omega_n^2},\qquad
M_p=e^{-\pi\zeta/\sqrt{1-\zeta^2}}\ (0<\zeta<1),\quad
t_s\approx\frac4{\zeta\omega_n}.
\end{equation}
\begin{equation}
\label{eq:settling-band}
|e^{-t/\tau}|\leq \varepsilon
\quad\Rightarrow\quad
t_s(\varepsilon)\approx -\tau\ln(\varepsilon).
\end{equation}
For \(2\%\) giver \(-\ln(0.02)\approx3.91\), som afrundes til \(4\tau\). For \(1\%\)
giver \(-\ln(0.01)\approx4.61\), så et langsomt førsteordensled med
\(\tau=5\) s har \(t_s\approx23\) s, dvs. nærmeste eksamenssvar kan være 25 s.

\paragraph{Symboler}
\(\tau\) [s] er tidskonstant; \(\omega_n\) [rad/s] naturlig frekvens;
\(\zeta\) dimensionsløs dæmpning; \(M_p\) er overshoot som ratio.

\paragraph{Antagelser}
Overshootformlen gælder standard underdæmpet anden orden. Approksimationen for
\(t_s\) er en 2\%-regel.

\paragraph{Hurtigtjek}
Et stabilt førsteordenssystem overshooter ikke. Større \(\zeta\) betyder mindre
overshoot; højere \(\omega_n\) betyder typisk hurtigere respons.
Closed-loop bandwidth fra et Bodeplot aflæses ved første frekvens hvor magnitude er
3 dB under lavfrekvensniveauet:
\[
|T(j\omega_b)|=\frac{|T(0)|}{\sqrt2}.
\]
Hvis lavfrekvensniveauet er \(+3\) dB, ligger \(-3\) dB-punktet derfor ved 0 dB i
plottet. For et oscillerende steprespons aflæses den dæmpede egenfrekvens fra
peakperioden \(T_d\):
\[
\omega_d=\frac{2\pi}{T_d}.
\]

\paragraph{Typiske fælder}
At bruge procent i formlen uden at dividere med 100; F21 Q9 afrunder grænsen
\(K=20\), som giver \(12.026\%\) ved eksakt beregning. Ved match mellem Bode og
steprespons skal både slutværdi, initialhop, resonans og hastighed passe; én af dem er
ikke nok.

\paragraph{Python}
\code{second\_order\_characteristics(denominator)} returnerer
\(\omega_n,\zeta,M_p\) for en allerede identificeret standardnævner og kan kontrollere
F21 Q9; vurder selv om højereordensdynamik gør metoden ugyldig.
\code{second\_order\_analysis(denominator)} bruges, når du også vil have poler,
dæmpningsklasse, dæmpet frekvens, peak time og 2\%/5\%-settling-estimater. Den kræver
stadig, at nævneren faktisk må behandles som standard anden orden.

\subsection{Dominerende poler og delay / model reduction}
\label{sec:dominant-delay}

\paragraph{Brug når}
Slides spørger efter mest dominerende pol, om ekstra poler kan ignoreres, eller når et
procesrespons har tydelig time delay/lag.

\paragraph{Metode}
\begin{enumerate}
\item Dominerende poler er de stabile poler tættest på imaginæraksen, dvs. med mindst
      negativ realdel. De bestemmer den langsomste envelope i tidsresponsen.
\item For en reel pol \(p<0\) er tidskonstanten \(\tau=-1/p\). For et komplekst par
      \(-\zeta\omega_n\pm j\omega_d\) er envelope-tidskonstanten
      \(1/(\zeta\omega_n)\), og oscillationsperioden er \(2\pi/\omega_d\).
\item En hurtig ekstra pol kan ofte ignoreres i en dominant-pole-approksimation, hvis
      dens realdel ligger mindst ca. 5 gange længere mod venstre end den dominerende
      pol. Ellers kan den ændre overshoot, rise time og Bodefase mærkbart.
\item Time delay \(L\) skrives ideelt som \(e^{-sL}\). Det ændrer ikke magnitude, men
      giver ekstra fase \(-\omega L\) rad, altså \(-\omega L\,180/\pi\) grader.
\item Hvis opgaven kræver en rational model for delay, bruges en Padé-approksimation
      kun som approximation; den kan introducere zeros og må ikke forveksles med den
      fysiske delay.
\end{enumerate}

\paragraph{Hurtigtjek}
I controllerdesign skal \(\omega_c\) ikke vælges så højt, at delayfasen eller en
ikke-dominerende pol spiser phase margin. Ved limited systems-slides vælges
\(\omega_c\) ofte over den dominerende pol for hastighed, men ikke så højt at rate
limiteren aktiveres.

\paragraph{Typiske fælder}
At vælge polen med størst absolutværdi som dominerende; at ignorere delayfase ved
crossover; at bruge en dominant-pole-approksimation uden at tjekke de øvrige poler.
